\maketitle

\begin{abstract}
This report reconstructs as much as possible of the mathematical content of the
Wikipedia article \emph{Expected value} inside intuitionistic primitive
recursive arithmetic, abbreviated \(\ipra\).  The central observation is that a
large majority of the article's algebraic probability theory is already
finitary.  On a finite sample space with natural weights and rational-valued
observables, expectation is one bounded sum divided by a positive integer.
Normalization, linearity, monotonicity, the law of the unconscious statistician,
indicator and tail formulas, variance identities, Markov and Chebyshev
inequalities, Cauchy--Schwarz, independence, finite conditional expectation,
and sample-mean formulas reduce to primitive-recursive calculations and
quantifier-free induction.

The article also discusses countable laws, densities, general Lebesgue
integration, convergence theorems, characteristic functions, and applications.
Those statements are not uniformly meaningful in a theory whose only objects
are natural numbers.  We give precise effective replacements: every infinite
object is a primitive-recursive stream of finite rational approximants together
with an explicit tail, Cauchy, integrability, or continuity modulus.  The
resulting claims are theorem schemata of \(\ipra\), not appeals to an ambient
classical universe.  We isolate where a classical converse---most notably an
existence-only Radon--Nikodym assertion---would require additional data and is
therefore not claimed.

The manuscript supplies complete human-readable arguments, a theorem-by-theorem
coverage ledger, three Lean design files, a route to explicit derivations in
ProveIt's \code{Prf PRAAx}, and an exact-rational regression program.  The Lean
files are blueprints rather than a completed kernel-checked formalization.  The
Python run checks 148,317 exact assertions and uses no floating-point value as a
proof oracle.
\end{abstract}

\vspace{1em}
\noindent\textbf{Formal status.}
The mathematical proofs below are written so that every finite step can be
compiled to bounded sums, products, comparisons, pairing, and
quantifier-free induction.  The effective analytic arguments state all moduli
needed for such a compilation.  Nevertheless, this packet does not contain the
final explicit \code{PRAAx} proof terms.  Statements marked ``internal'' are
intended iPRA theorems; statements marked ``effective analytic'' are iPRA
schemata for a fixed primitive-recursive presentation and its certificates.

\tableofcontents
\newpage

\section*{Executive theorem map}
\addcontentsline{toc}{section}{Executive theorem map}

The paper has three layers.

\begin{enumerate}
\item \textbf{Exact finite layer.}  \Cref{sec:coding,sec:finite,sec:properties,sec:inequalities,sec:independence,sec:samples}
  develops expectation on a finite weighted law.  No real numbers, completed
  infinite sets, or measure spaces are used.  After clearing denominators, the
  statements are equations and inequalities over \(\N\).
\item \textbf{Effective limiting layer.}  \Cref{sec:countable,sec:integration,sec:distributions,sec:convergence}
  treats series and integrals through finite approximants and explicit moduli.
  The finite kernel is applied at a precision computed from the requested output
  precision.
\item \textbf{Formalization layer.}  \Cref{sec:fourier,sec:lean,sec:conclusion}
  covers finite Fourier methods and applications, records the Lean interfaces,
  and identifies the remaining proof-engineering work.
\end{enumerate}

\begin{center}
\begin{tabular}{@{}p{0.24\textwidth}p{0.66\textwidth}@{}}
\toprule
Class & Representative results \\
\midrule
Exact in \(\ipra\) & finite expectation and LOTUS; linearity; monotonicity;
positive/negative parts; indicator and tail formulas; variance; Markov,
Chebyshev and Cauchy--Schwarz; product-law independence; finite conditioning;
sample mean; finite Kolmogorov inequality. \\
Effective schema & absolutely convergent countable expectation; effective
reordering; lower-real extended expectation; compact and improper integration;
monotone/Fatou/uniform-integrability convergence; common continuous laws;
characteristic-function differentiation. \\
Constructive correction & conditionally convergent sums are order sensitive;
Cauchy truncations are path dependent; pointwise convergence does not imply
convergence of expectations; the article's ``about 53 percent'' die statement
corresponds to an open interval, not the inclusive event. \\
Not claimed uniformly & extraction of a density from bare absolute continuity;
unwitnessed arbitrary Lebesgue objects; qualitative almost-sure convergence
without a finite-error contract or a modulus presentation. \\
\bottomrule
\end{tabular}
\end{center}
